\section{Euclid's engine: the impossibility of perpetual descent}\label{sec:engine}

The programme reduces the twin-prime conjecture to the behaviour of one dynamical object --- a
chain of self-similar Euclidean equations we call \emph{Euclid's engine} --- and rests the whole
edifice on a single fact: the engine cannot run forever. Suppose there were only finitely many
twin primes. Then, from some scale onward, every pair $(6m-1,\,6m+1)$ has a composite side whose
factorisation exposes an active divisor; removing it yields a structurally identical Euclidean
equation with a strictly smaller centre. Iterating, one obtains an infinite chain of clean
descents that never leaves the clean core --- an eternally running engine performing the
nontrivial work of descent, losing nothing. We prove that no such engine exists. This is Fermat's
infinite descent, restated for the centres of Euclidean pairs; the state is described by a single
natural number --- its \emph{height} --- and one successful step decreases the height strictly and
\emph{multiplicatively}, after which only the arithmetic of $\N$ does the work. Everything below is
machine-checked on the bare Lean~4 kernel, without \code{mathlib}
(\code{EuclidsPath/Engine/EPMI.lean}, \code{EuclidsPath/Engine/Irreversibility.lean},
\code{EuclidsPath/Engine/UniversalEngine.lean})~\cite{lean4,euclidspath-repo}.

\subsection{The descent step and its strictness}

We strip away the arithmetic filling and keep only what the impossibility proof needs. The state
carries one numerical characteristic, the height $h\in\N$, substantively the centre $m$ of the
pair $(6m-1,\,6m+1)$.

\begin{definition}[\code{DescentStep}]\label{thm:DescentStep}
Let $A\in\N$ be a fixed scale threshold (a lower bound on the active factors). One successful clean
descent from height $h$ to height $h'$ is the relation
\begin{equation}\label{eq:descentstep}
  \mathrm{DescentStep}\,A\,h\,h' \;:=\; A\cdot h' < h.
\end{equation}
\end{definition}

The step fixes not merely a decrease but a \emph{multiplicative} one: the new height is smaller by
a factor of $A$ (up to strictness). It is the factor $A$, not the difference, that sets the speed
of descent and makes termination inevitable.

\begin{lemma}[\code{descent_strict}]\label{thm:descent_strict}
\green{}\; If $1\le A$ and $\mathrm{DescentStep}\,A\,h\,h'$, then $h'<h$.
\end{lemma}

\emph{Proof idea.} From $1\le A$ we get $h'\le A\cdot h'$ (multiplying by a positive factor does
not decrease), and $A\cdot h'<h$ by \cref{eq:descentstep}; transitivity gives $h'<h$. The lemma
separates the \emph{direction} of the descent (the height falls) from its \emph{speed} (it falls by
a factor of $A$).

\subsection{The impossibility of infinite descent}

\begin{theorem}[\code{no_infinite_descent}]\label{thm:no_infinite_descent}
\green{}\; Let $1\le A$. There is no sequence of heights $H:\N\to\N$ for which every step is a
successful $A$-descent, i.e.
\begin{equation}\label{eq:perpetualheights}
  \forall\, t,\quad \mathrm{DescentStep}\,A\,(H\,t)\,(H\,(t+1)).
\end{equation}
From the existence of such a sequence, \code{False} is derived.
\end{theorem}

\begin{proof}
Consider the quantity $H\,t+t$. At each step the height falls by at least $1$ (\cref{thm:descent_strict}),
while the counter $t$ grows by $1$; hence $H\,t+t$ does not increase, and by induction
$\forall\,t,\ H\,t+t\le H\,0$. Substituting $t=H\,0+1$ gives
$H\,(H\,0+1)+(H\,0+1)\le H\,0$, whose left side already exceeds $H\,0$ --- a contradiction. In Lean
the induction is closed by \code{omega}; \code{\#print axioms} confirms the theorem depends on no
axioms whatsoever and is free of \code{sorry}.
\end{proof}

The multiplicative decrease $H_t<H_0/A^{t}$ drives the height below $1$ in finitely many steps, and
a positive integer cannot be less than $1$: the natural numbers have a bottom, and a discrete
descent must reach it. The contradiction needs only $A=1$ (the well-orderedness of $\N$); the
factor $A\ge 1$ supplies the quantitative speed. This is stronger than the physical second law:
because the height is an integer with a strict bottom, the descent does not merely die down
asymptotically --- it halts in at most $H\,0$ steps (\cref{thm:turned_engine_halts}).

\subsection{The structural form: state and boundary-exit}

To connect the abstract theorem to the substantive picture we introduce an explicit state and a
step operator that distinguishes the two outcomes of removing an active factor.

\begin{definition}[\code{State}]\label{thm:State}
The engine state is the structure with a single field, $\mathrm{height}:\N$; substantively this is
the centre $m$ of the current Euclidean pair. A step over a state $s$ is an inductive type with two
constructors,
\begin{equation}\label{eq:step}
  \mathrm{Step}\,A\,s \;=\;
  \begin{cases}
    \texttt{clean}\ s'\ (A\cdot s'.\mathrm{height} < s.\mathrm{height}), & \text{successful descent into a clean state,}\\[2pt]
    \texttt{boundary}, & \text{the absorbing exit } \bot,
  \end{cases}
\end{equation}
where the \texttt{clean} branch carries a witness of the multiplicative decrease, and
\texttt{boundary} records exit from the clean core onto the boundary, from which there is no
return.
\end{definition}

\begin{theorem}[\code{no_perpetual_engine}]\label{thm:no_perpetual_engine}
\green{}\; Let $1\le A$. There is no trajectory $\mathrm{run}:\N\to\mathrm{State}$ in which every
step is a successful clean descent, i.e.
$A\cdot(\mathrm{run}\,(t+1)).\mathrm{height} < (\mathrm{run}\,t).\mathrm{height}$ for all $t$.
\end{theorem}

\emph{Proof idea.} Reduce to the abstract form: take the heights $t\mapsto(\mathrm{run}\,t).\mathrm{height}$
and apply \cref{thm:no_infinite_descent}. This is precisely ``there is no perpetual engine of
Euclid''.

\begin{theorem}[\code{boundary_dichotomy}]\label{thm:boundary_dichotomy}
\green{}\; For any step $st:\mathrm{Step}\,A\,s$, exactly one of two things holds:
\begin{equation}\label{eq:dichotomy}
  \bigl(\exists\, s'\,h,\ st = \texttt{clean}\ s'\ h\bigr)\ \lor\ \bigl(st = \texttt{boundary}\bigr).
\end{equation}
\end{theorem}

\emph{Proof idea.} Case analysis over the constructors of \cref{eq:step}. This is the type-level
fixation of the boundary-exit law: a descended centre either stays clean (height strictly fell by a
factor of $A$) or exits onto the absorbing boundary $\bot$, from which no clean continuation of the
same branch exists. The boundary absorbs: $\bot\not\to S$.

\subsection{Irreversibility and the arrow of time}

\cref{thm:no_infinite_descent} answers ``does the engine always halt?''. The companion module
\code{Irreversibility.lean} closes ``can it turn back?'', supplying a discrete analogue of the
second law of thermodynamics. All of it is machine-proven on the standard axioms alone, with no
appeal to the distribution of primes.

\begin{theorem}[\code{engine_never_returns}]\label{thm:engine_never_returns}
\green{}\; Let $1\le A$ and let $H:\N\to\N$ satisfy \cref{eq:perpetualheights}. Then $H$ is
strictly antitone:
\begin{equation}\label{eq:strictanti}
  \mathrm{StrictAnti}\,H \;:\Longleftrightarrow\; \bigl(\forall\, s\,t,\; s<t \Rightarrow H\,t<H\,s\bigr).
\end{equation}
\end{theorem}

\emph{Proof idea.} The local strict decrease $H(n+1)<H(n)$ from \cref{thm:descent_strict} is lifted
to global antitonicity by \code{strictAnti_nat_of_succ_lt}: on $\N$, decrease at every unit step is
equivalent to decrease over any interval. The height is the coordinate of thermodynamic time, and
\cref{eq:strictanti} is the formal arrow of time --- the engine, once descending, cannot restore a
previous state. This describes any chain of successful descents \emph{should one be given}; whether
such a chain can be infinite is settled negatively by \cref{thm:no_infinite_descent}.

\begin{theorem}[\code{fuel_ascent_strictMono}]\label{thm:fuel_ascent_strictMono}
\green{}\; The map $n\mapsto n+1$ is strictly monotone: $\mathrm{StrictMono}\,(\lambda n.\,n+1)$.
\end{theorem}

\emph{Proof idea.} Trivially $n<n+1$. The triviality of form hides the substance: the successor
chain $0,1,2,\dots$ is an infinite strictly increasing trajectory. Adding $+1$ is ``fuel'': larger
centres never run out. Setting this beside \cref{thm:no_infinite_descent} exhibits the asymmetry of
the thermodynamic axis --- infinitely up, finitely down. That asymmetry is the arrow.

\begin{theorem}[\code{turned_engine_halts}]\label{thm:turned_engine_halts}
\green{}\; Let $H:\N\to\N$ and suppose the engine has made $k$ strict steps downward, i.e.
$H(t+1)<H(t)$ for all $t<k$. Then
\begin{equation}\label{eq:haltbound}
  k \;\le\; H(0).
\end{equation}
\end{theorem}

\emph{Proof idea.} By induction the invariant $H(t)+t\le H(0)$ holds for all $t\le k$; substituting
$t=k$ and dropping $H(k)\ge 0$ gives \cref{eq:haltbound}. This is the same ``height plus time does
not grow'' balance as in \cref{thm:no_infinite_descent}, read as a finite bound rather than a route
to contradiction: any turn downward is a finite path of length at most the initial height.

\subsection{The universal form: engine versus well-foundedness}

The impossibility above is a shadow of a single order-theoretic fact, formalised in
\code{UniversalEngine.lean} over an arbitrary relation. This is the reading the whole article
carries: where a carrier is well-founded, the engine is forbidden \emph{unconditionally}; where it
is not, the engine operates.

\begin{definition}[\code{PerpetualEngine}]\label{thm:PerpetualEngine}
For a relation $r$ on a type $\alpha$, a perpetual engine is an infinite strictly $r$-descending
chain:
\begin{equation}\label{eq:perpetualengine}
  \mathrm{PerpetualEngine}\,r \;:=\; \exists\, f:\N\to\alpha,\ \forall\, n,\ r\,(f(n+1))\,(f\,n).
\end{equation}
\end{definition}

\begin{theorem}[\code{perpetualEngine_iff_not_wellFounded}]\label{thm:perpetualEngine_iff_not_wellFounded}
\green{}\; For any $r$ on $\alpha$,
\begin{equation}\label{eq:dividingline}
  \mathrm{PerpetualEngine}\,r \;\Longleftrightarrow\; \neg\,\mathrm{WellFounded}\,r.
\end{equation}
\end{theorem}

\emph{Proof idea.} Unfolding $\mathrm{WellFounded}$ via its descending-chain characterisation, an
engine \cref{eq:perpetualengine} is literally a descending chain, whose non-existence is
well-foundedness. This is the exact dividing line: the engine and well-foundedness are complementary.

\begin{theorem}[\code{no_perpetual_engine_of_wellFounded}]\label{thm:no_perpetual_engine_of_wellFounded}
\green{}\; If $\mathrm{WellFounded}\,r$, then $\neg\,\mathrm{PerpetualEngine}\,r$.
\end{theorem}

\emph{Proof idea.} Immediate from \cref{thm:perpetualEngine_iff_not_wellFounded}. The impossibility
is intrinsic, almost definitional: the engine's own defining chain is exactly what
well-foundedness forbids. \cref{thm:no_infinite_descent} is the instance $r={<}$ on $\N$
(\code{no_perpetual_engine_on_nat}), with height as rank.

\begin{theorem}[\code{no_perpetual_engine_of_rank}]\label{thm:no_perpetual_engine_of_rank}
\green{}\; Let $\rho:\alpha\to\beta$ with $(\beta,s)$ well-founded, and suppose every $r$-step
strictly decreases the rank: $\forall\, x\,y,\ r\,x\,y \Rightarrow s\,(\rho\,x)\,(\rho\,y)$. Then
$\neg\,\mathrm{PerpetualEngine}\,r$.
\end{theorem}

\emph{Proof idea.} An engine $f$ for $r$ would transport along $\rho$ to a descending chain
$\rho\circ f$ in the well-founded $(\beta,s)$, contradicting
\cref{thm:no_perpetual_engine_of_wellFounded}. The $\N$-rank specialisation
(\code{no_perpetual_engine_of_natRank}) is the unification: every descent front of the programme is
a corollary of this one theorem, obtained by exhibiting a rank that strictly decreases along steps.

\begin{theorem}[\code{perpetualEngine_on_real}]\label{thm:perpetualEngine_on_real}
\green{}\; On the continuum the engine operates:
$\mathrm{PerpetualEngine}\,(<)$ on $\R$.
\end{theorem}

\emph{Proof idea.} The witness $(1/2)^n$ is an infinite strictly $(<)$-descending chain in $\R$;
hence, by \cref{thm:perpetualEngine_iff_not_wellFounded}, the strict order on $\R$ is not
well-founded. This marks the honest controller boundary: on discrete carriers the engine is
forbidden and descent halts, while on the reals perpetual descent is realised. The prohibition
lives in discreteness, not in analysis.

\medskip

We have introduced the central object and proved, machine-checked and axiom-free, that it cannot
run forever (\cref{thm:no_infinite_descent,thm:no_perpetual_engine}), does not turn back
(\cref{thm:engine_never_returns}), and, once descending, halts in finitely many steps
(\cref{thm:turned_engine_halts}); the dichotomy (\cref{thm:boundary_dichotomy}) leaves the absorbing
boundary as the only alternative to continuing. We draw the boundary of what is proved honestly:
this is the impossibility of an infinite \emph{local} clean descent along a single branch. The
global non-covering result does \emph{not} follow automatically --- one can imagine a finite tree of
bounded depth, each branch honestly terminating in an absorbing boundary-leaf, that nevertheless
covers all centres. The closure plan lifts the prohibition from a branch to the whole tree, recast
as a multi-rank non-cover with a rank reduction to the rank-$1$ neighbour obstruction; that single
open node is taken apart in \cref{sec:reduction}. The engine, for its part, supplies the local law
of finiteness on which the reduction stands.
